%قبل از شروع مسئله ترجیحاً چیزی ننویسید. اما اگر خواستید بنویسید، در قالب لیستی (فقرات و شمارش) بنویسید.
%برای اضافه کردن نمره، بعدی از شروع‌مسئله، [x] را اضافه کنید و به جای x  نمره مورد نظر را قرار دهید. برای زیرمسئله‌ها هم به همین صورت است.
%\شروع‌مسئله[10]
\شروع‌مسئله[25]
سیل‌های نوروزی بسیاری از راه‌های ارتباطی را مسدود کرد. در این مسئله $n$ روستا داریم که راه‌های ارتباطی آن فاقد دور است. در هر مرحله با بالا آمدن سیل، یکی از راه‌های ارتباطی قطع می‌شود. در هر مرحله شما باید اندازهٔ بزرگترین زیرمجموعه از روستاها را محاسبه کنید که به هم هنوز مسیر دارند. 

در خط اول ورودی، تعداد روستا‌ها $n$ ($1\leq n\leq 1,000,000$) و تعداد راه‌های ارتباطی آن‌ها $m$ داده شده است. سپس در $m$ خط بعدی، در هر خط دو عدد $u$ و $v$ آمده است که $u$ و $v$ به هم راه دوطرفه دارند. روستا‌ها از $1$ تا $n$ شماره‌گذاری شده‌اند. سپس در خط بعدی عدد $q$ آمده و در $q$ خط بعدی شمارۀ راه‌ارتباطی‌ای که قطع می‌شود آمده است. راه‌ها به ترتیبی که در ورودی آمده‌اند شماره‌گذاری می‌شوند. خروجی باید $q$ خط داشته باشد که در هر خط اندازهٔ بزرگترین زیرمجموعه از روستاها که پس از قطع شدن راه ارتباطی آن مرحله هنوز به هم مسیر دارند آمده باشد.

راهنمایی: توجه کنید که کل ورودی از ابتدا به شما داده شده است و نیاز نیست سؤال‌ها را به ترتیب جواب دهید. در این سؤال شما باید داده‌ساختار مجموعه‌های مجزا که در کلاس مطرح شد را پیاده‌سازی و از آن استفاده کنید.

\begin{center}
	\begin{tabular}{|c|c|}
		\hline
		خروجی                 
		& ورودی    \\ \hline
\parbox[t]{5cm}{\normalsize\raggedright{\texttt{3\\2}}\vspace{0.2cm}} &
\parbox[t]{5cm}{\normalsize\raggedright{\lr{\texttt{5 4\\1 2\\2 3\\1 4\\1 5\\2\\1\\4}}}\vspace{0.2cm}}\\\hline
			\end{tabular}
	\end{center}
%\زیرمسئله % مسئلۀ ۱پ
%\زیرمسئله % مسئلۀ ۱ت


\solution{
با توجه به اینکه کلیه اطلاعات مسئله از ابتدا در اختیار برنامه قرار دارد، می‌توان اتفاقات را از انتها به ابتدا پاسخ داد. در این صورت در ابتدا یک گراف احتمالاً ناهمبند داریم که در هر مرحله یکی این که اندازۀ بزرگترین مؤلفۀ همبندی را باید خروجی دهید و دوم این‌که یک یال به گراف اضافه می‌شود. اضافه کردن یال با داده ساختار مجموعه‌های مجزا به راحتی قابل انجام است و به روز رسانی اندازۀ بزرگترین مؤلفۀ همبندی با نگه‌داری size (مثل \متن‌لاتین{rank}) برای نماینده هر مجموعه و به‌ روز رسانی آن در زمان ادغام دو مجموعه قابل انجام است.

کد پایتون این سؤال ضمیمه شده است.

}




%%%%%%%%%%%%%%%%%%%%%%%%%%%%%%%%%%%%%%%%%%%%%%%%%%%%%
% برخی دستورات مفید برای تمرین‌ها. %
%%%%%%%%%%%%%%%%%%%%%%%%%%%%%%%%%%%%%%%%%%%%%%%%%%%%%

% فرمول‌های نمونه:
%$f(x) = \Theta(x)$
%$T(x, y) \leq \log(x) + 2^y + \binom{2n}{n}$
% {\tt code\_function}


% لیست موارد بدون شماره‌گذاری:
%\شروع{فقرات}
%    \فقره مورد اول 
%        \شروع{فقرات}
%            \فقره مورد اول داخلی 
%                \شروع{فقرات}
%                    \فقره اولین مورد خیلی داخلی 
%                    \فقره دومین مورد خیلی داخلی 
%                \پایان{فقرات}
%            \فقره مورد دوم داخلی 
%        \پایان{فقرات}
%    \فقره مورد دوم 
%\پایان{فقرات}

% لیست موارد شماره‌گذاری شده:
%\شروع{شمارش}
%    \فقره یک مورد 
%    \فقره یک مورد دیگر 
%    \فقره باز هم یک مورد دیگر
%\پایان{شمارش}

% برای نوشتن یک فرمول چند خطی که یک شماره بخوره. برای حالت‌های دیگر به سایت زیر مراجعه کنید:
% https://www.overleaf.com/learn/latex/Aligning_equations_with_amsmath
%\begin{align} 
%    \begin{split}
%        (x+y)^3 &= (x+y)^2(x+y) \\
%                &= (x^2+2xy+y^2)(x+y) \\
%                &= (x^3+2x^2y+xy^2) + (x^2y+2xy^2+y^3) \\
%                &= x^3+3x^2y+3xy^2+y^3
%    \end{split}                                 
%\end{align}

% نوشتن ماتریس و آرایه دوبعدی
%\begin{align}
%    A = 
%    \begin{bmatrix}
%        A_{11} & A_{21} \\
%        A_{21} & A_{22}
%    \end{bmatrix}
%\end{align}

% اضافه کردن تصویر با شرح. برای اضافه کردن تصویر بدون شرح، فقط دستور درج‌تصویر را استفاده کنید.
%\شروع{شکل}[h!]
%\تنظیم‌ازوسط
%  \درج‌تصویر[پهنا=0.5\پهنای‌متن]{img.jpg}
%\شرح{شرح تصویر}
%\برچسب{تصویر:۱}
%\پایان{شکل}
%\پایان{زیرمسائل}

%کد زیر را ببینید.
%%\lstinputlisting{source_filename.py}
%\begin{latin}
%\begin{lstlisting}
%# Import the modules
%import sys
%import random
%
%ans = True
%
%while ans:
%question = raw_input("Ask the magic 8 ball a question: (press enter to quit) ")
%
%answers = random.randint(1,8)
%
%if question == "":
%sys.exit()
%
%elif answers == 1:
%print "It is certain"
%
%elif answers == 2:
%print "Outlook good"
%
%elif answers == 3:
%print "You may rely on it"
%
%elif answers == 4:
%print "Ask again later"
%
%elif answers == 5:
%print "Concentrate and ask again"
%
%elif answers == 6:
%print "Reply hazy, try again"
%
%elif answers == 7:
%print "My reply is no"
%
%elif answers == 8:
%print "My sources say no"
%\end{lstlisting}
%\end{latin}
